\subsection{Excluding order $44$}\label{sec:d31a44}
\begin{proposition}\label{prop:decic44-exclusion}
An elliptic curve over a number field of degree $10$ has no point of order $44$.
\end{proposition}
\begin{proof}
Put $X=X_1(44)$ and $J=J_1(44)$.
The genus is $36$, and $J(\mathbb Q)$ has rank zero by
\cite[Theorem~3.1]{DEHMZ}.
Moreover, $\operatorname{gon}_{\mathbb Q}(X)=15$ by
\cite[Theorem~D]{Varivoda26}.
Suppose $D$ is a noncuspidal effective rational divisor of degree $10$ on $X$.

We first exclude entirely cuspidal reduction at $3$.
Such a reduction lifts to an effective cuspidal divisor $C$ over $\mathbb Q_3$,
since the cyclotomic cusp fields are unramified at $3$.
Both $[D-10c_0]$ and $[C-10c_0]$ are torsion, by rank zero and
Manin--Drinfeld, respectively, where $c_0$ is a rational cusp.
Thus $[D-C]\in J(\mathbb Q_3)$ is torsion and has zero reduction.
The torsion kernel of reduction over $\mathbb Q_3$ is zero, giving $D\sim C$.
Since $D\ne C$, we have $h^0(X_{\mathbb Q_3},D)\ge2$.
The divisor $D$ is defined over $\mathbb Q$, so cohomology and base change give
$h^0(X,D)\ge2$, contradicting rational gonality $15$.
This argument does not require $C$ to be defined over $\mathbb Q$.

The noncuspidal closed-point counts at $3$ through degree $10$ are
\[
 (0,0,0,0,60,160,240,920,2260,5854).
\]
They are checked independently from the exact rational matrices of $T_3$
and $\langle3\rangle$ using Eichler--Shimura.
Consequently the noncuspidal part of $\overline D$ is either a single
place of degree $5\le d\le10$, or a pair of quintic places, with repetitions allowed.

The global cusp signature is $1^{15}2^5 5\,10^2$, whereas the signature
at $3$ is $1^{15}2^5 5^3 10$.
Every cusp of degree at most $4$ is globally intact.
The retained Hecke tests for singleton degrees $6$, $7$ and $8$ therefore
apply with the larger complements of degrees $4$, $3$ and $2$.
There are respectively $16$, $24$ and $92$ diamond representatives,
and all tests are nonprincipal.
The $226$ retained degree-$9$ canonical tests also apply, since their
new complement is a rational cusp annihilated literally by $A_5$.

For degree $5$, two quintic cusp orbits are fragments of a global
orbit of degree $10$. We use the cusp formula directly.
Write $\sigma_q$ for the cyclotomic Galois action at a good prime $q$.
On cuspidal divisors,
\[
 T_q=\sigma_q+q\langle q\rangle\sigma_q^{-1}.
\]
Since $5\equiv3^8\pmod {44}$, every Frobenius-$3$ cuspidal trace is fixed by
$\sigma_5$. Thus
\[
 A_5C'=0,\qquad A_5=T_5-1-5\langle5\rangle,
\]
as an actual divisor, including the split quintic traces.
The same identity holds for a rational cusp $c_0$.

For any good odd $q$, the operator $A_q=T_q-1-q\langle q\rangle$
annihilates $J(\mathbb Q)$ by Eichler--Shimura and injectivity of
rational torsion reduction at $q$.
Hence the following are necessary principal-divisor tests:
$A_5P_5$ for a singleton quintic noncuspidal part, and $A_qD'$ for
an entirely noncuspidal degree-$10$ reduction $D'$.
The complete new test list consists of
\[
\begin{array}{c|r}
\text{support type}&\text{tests}\\\hline
(5)&6\\
(10)&590\\
(5,5)&360
\end{array}
\]
The $590$ degree-$10$ diamond orbits cover all $5854$ places.
For $(5,5)$ we take one of the six quintic representatives for the first
factor and every one of the $60$ quintic places for the second.
This includes all pairs and repeated places.

We use the retained canonical model and multiplication tensor
\[
 R_2\otimes R_2\longrightarrow R_4,
 \qquad R_i=H^0(X_{\mathbb F_3},iK_X),\qquad
 \dim R_2=105,\quad\dim R_4=245.
\]
The canonical coordinates and all $11025$ products are certified on the
function-field model. Since $\deg(2K_X)=140\ge2g+1$, this identifies the
computed products with the full space $R_4$.
The old degree-$9$ multiplication certificates use this same model and tensor.

For a new support $D'$ put $T=T_qD'$ and $V=D'+q\langle q\rangle D'$.
The actual marked cyclic isogenies give $T$, including residue-field norms.
We merge coincident closed places on each side before computing local jets.
Thus the tests also cover the stabilizers that can occur in even degree.
The two divisors have degree $30$ for the singleton quintic $A_5$ test,
$60$ for a degree-$10$ $A_5$ test, and $80$ for a degree-$10$ $A_7$ test.

Choose nonzero $F,H\in R_2(-T)$.
If $f\in\mathcal L(T-V)$, then $G=fF\in R_2(-V)$ and $GH\in FR_2$.
It is therefore enough to show that
\[
 R_2\longrightarrow H^0(V,(2K_X)|_V)\oplus R_4/FR_2,
 \qquad G\longmapsto(G|_V,GH\bmod FR_2)
\]
is injective.
Multiplication by $F$ has rank $105$, so the second summand has dimension $140$.
The exact restriction and multiplication matrices give rank $105$
for $954$ of the $956$ new tests with $q=5$.
The two remaining cases are singleton degree-$10$ supports.
Using $q=7$, both resulting matrices have rank $105$.
Thus every necessary new divisor is nonprincipal.
Together with the unchanged lower-degree tests, this excludes every
possible reduction of $D$.
\end{proof}

The new source, marked-isogeny and canonical-jet data are retained in
\path{continuation_20261031a/C44}.
The source count is independently checked from the rational Frobenius matrices.
The canonical arithmetic and marked-isogeny controls are replayed separately;
the two nonrejections for $A_5$ are retained and are resolved only by $A_7$.
